% ---------------------------------------------------------------------------
% PROSE-APEX deployment-mode boundary -- MAIN figure (1x4).
% Generated by experiments/run_mode_boundary.py; numbers come from SimCXL.
% IMPORTANT: width=\linewidth fills the full column -- without it the panels
% render small with side margins. Use \columnwidth for a single-column float,
% or wrap in figure* + \textwidth for a full-page-width float in two-column.
% ---------------------------------------------------------------------------
\begin{figure}[t]
  \centering
  \includegraphics[width=\linewidth]{figures/fig_mode_boundary_main.pdf}
  \caption{\textbf{Deployment-mode boundary and robustness.}
    Mode~A is the endpoint-push fast path (peer-to-peer DMA, performance upper
    bound). Mode~B uses endpoint-gated pull: it preserves \emph{zero} RPE while
    paying a small extra pull latency per batch. Mode~C is a passive Type-3
    software fallback that is single-host safe but reopens a multi-host
    atomicity gap.
    \textbf{(a)}~Under candidate oversubscription on a starved link, A/B deliver
    the full speedup over FullKV while C tracks close only until the link
    saturates.
    \textbf{(b)}~A and B are correctness-equivalent (both bind admission at the
    endpoint) and hold RPE${=}0$ at every host count; C's residual RPE reopens
    toward the reported 16-host bound of $14.4\%$.
    \textbf{(c)}~Mode-B pull-path sensitivity: the cost of dropping endpoint DMA
    is $2$--$5\,\mu$s/batch, a single-digit throughput loss relative to Mode~A,
    while RPE stays $0$; host-bounce pull is strictly worse than P2P pull.
    \textbf{(d)}~The $14.4\%$ residual is \emph{not} a single-trace artifact:
    across workload and policy variants it persists, decomposed into an
    eviction-race component (moves with workload/policy stress) and a
    protocol-driven epoch-rollover component (roughly flat). Only the
    endpoint-gated A/B path eliminates RPE for \emph{every} setting.}
  \label{fig:mode-boundary}
\end{figure}
